\documentclass[conference]{IEEEtran}
\usepackage[utf8]{inputenc}
\usepackage{amsmath,amssymb}
\usepackage{booktabs}
\usepackage{graphicx}
\usepackage{xcolor}
\usepackage[colorlinks=true,linkcolor=blue,citecolor=blue,urlcolor=blue]{hyperref}

% ============================================================================
% WORKSHOP PAPER 3 -- systems / edge-deployment + MEMORY BENCHMARKS
% Footprints measured from the live model node (load + inference); planner now ALSO
% measured on-device on a Jetson AGX Orin 64GB (JetPack 6, torch 2.8). Perception
% throughput still reasoned from device specs. Marked clearly where estimated.
% ============================================================================

\title{\textbf{The Planner Is Free, Perception Is the Wall:\\
Memory Benchmarks for Edge Deployment of Quantized Driving VLAs}}
\author{\IEEEauthorblockN{Justin Williams, Kishor Datta Gupta, Roy George, Mrinmoy Sarkar}
\IEEEauthorblockA{Clark Atlanta University \\ \texttt{justinwilliamstech693@gmail.com}}}

\begin{document}
\maketitle

\begin{abstract}
We characterize the memory and deployment profile of OpenDriveVLA-0.5B---a Qwen2.5-0.5B
language/action head on UniAD perception---toward on-board (NVIDIA Jetson Orin) inference. The
central, mundane-but-important finding: the language head everyone tries to quantize is already
trivial for the edge ($\approx 1$\,GB FP16, $\approx 0.25$\,GB at INT4; the whole inference node
fits in $3$--$5$\,GB), and \emph{group-wise} INT4 (the form Orin's INT4 kernels expect) adds only
$\approx 0.1$ bits/weight while preserving generation coherence. The wall is \emph{perception}:
UniAD's ResNet-101 over six $1600\times900$ cameras plus temporal BEVFormer dominates wall-clock,
runs FP32 with custom CUDA ops, and is far from real-time on Orin without a TensorRT-INT8 port. We
present in-memory benchmarks for the components and quantization levels---including \emph{on-device
measurements on a Jetson AGX Orin 64\,GB} (JetPack~6, torch~2.8): the 357.8\,M-parameter planner
loads and runs in $2.88$\,GB peak at $14.9$ tok/s, with a static weight footprint dropping
$0.72\to 0.19$\,GB ($3.9\times$) at group-wise INT4---and argue that edge VLA
research effort is mis-allocated: the planner is free; perception is the gating cost.
\end{abstract}
\begin{IEEEkeywords}
vision--language--action models, autonomous driving, post-training quantization, group-wise INT4, closed-loop evaluation, generation coherence, edge deployment, Jetson Orin
\end{IEEEkeywords}

\section{Introduction}
A recurring narrative is ``quantize the VLA language model and ship it to the robot.'' We provide
component-level memory benchmarks showing that, for a driving VLA, this optimizes the cheap part.
We quantify the planner footprint at FP16/INT8/INT4 (and the negligible group-wise overhead), the
total inference footprint, and the perception bottleneck, and map feasibility onto Jetson Orin and,
for contrast, a Raspberry Pi.

\section{Memory benchmarks}
\begin{table*}[t]\centering
\caption{Backbone (Qwen2.5-0.5B, 357.8\,M quantized Linear params) weight footprint and overhead.
Effective bits include per-group FP16 scales.}
\begin{tabular}{lccc}
\toprule
Precision & bits/weight (eff.) & backbone weights & vs FP16 \\
\midrule
FP16              & 16.0   & $\approx 1.0$\,GB  & $1.0\times$ \\
INT8 per-channel  & 8.06   & $\approx 0.50$\,GB & $0.50\times$ \\
INT4 per-channel  & 4.06   & $\approx 0.25$\,GB & $0.25\times$ \\
INT4 group-128    & 4.125  & $\approx 0.26$\,GB & $0.26\times$ \\
INT4 group-64     & 4.25   & $\approx 0.27$\,GB & $0.27\times$ \\
\bottomrule
\end{tabular}
\end{table*}
Group-wise INT4 costs $16/g$ extra bits/weight ($0.125$ at $g{=}128$): a $\sim$1\% footprint
increase over per-channel for the coherence robustness reported in our closed-loop study.

\begin{table*}[t]\centering
\caption{End-to-end inference footprint (measured on the live model node, single A100; the
language head is a small fraction). UniAD perception dominates compute and is FP32.}
\begin{tabular}{lcc}
\toprule
Component & GPU memory (load$\to$infer) & precision \\
\midrule
Qwen-0.5B head (FP16)        & $\approx 1$\,GB             & FP16 / INT4 \\
UniAD track+map perception   & bulk of $3$--$5$\,GB total  & FP32 (custom CUDA ops) \\
\textbf{Full node}           & \textbf{$\approx 2.5\to 5$\,GB} & mixed \\
\bottomrule
\end{tabular}
\end{table*}

\section{On-device measurement (Jetson AGX Orin 64\,GB)}
We ran the planner on a real \textbf{NVIDIA Jetson AGX Orin Developer Kit} (64\,GB, JetPack~6 /
L4T~R36.4.7, CUDA capability~8.7, PyTorch~2.8), loading the \emph{actual} OpenDriveVLA-0.5B
Qwen2.5 decoder weights (291/291 tensors) and decoding $60$ trajectory tokens from a $1024$-token
context (eager attention, default power mode, 30~iterations). Because our quantizer is fake-quant
(dequantized to FP16, no packed-INT4 GEMM in this build), latency and \emph{runtime} memory are
identical across precisions---confirming that the deployable win is the \emph{static weight
footprint}, not speed, until a packed kernel is added.
\begin{table*}[t]\centering
\caption{\textbf{Measured on Jetson AGX Orin 64\,GB} (JetPack 6, torch 2.8): OpenDriveVLA-0.5B
planner, 1024-token context $\to$ 60 generated tokens, median of 30 runs. Packed-weight column is
the true footprint if stored as packed INT$b$ + FP16 group scales.}
\begin{tabular}{lcccc}
\toprule
Config & median latency & throughput & runtime peak & packed weights \\
\midrule
FP16             & $4.02$\,s & $14.9$ tok/s & $2.88$\,GB & $0.716$\,GB \\
INT8 group-128   & $4.01$\,s & $15.0$ tok/s & $2.88$\,GB & $0.363$\,GB \\
INT4 group-128   & $4.02$\,s & $14.9$ tok/s & $2.88$\,GB & \textbf{$0.185$\,GB} \\
INT4 per-channel & $4.02$\,s & $14.9$ tok/s & $2.88$\,GB & $0.179$\,GB \\
\bottomrule
\end{tabular}
\end{table*}
The planner loads in $13.6$\,s and the whole component runs in under $3$\,GB on Orin---trivial for a
64\,GB board. Group-wise INT4 shrinks the stored planner $3.9\times$ ($0.716\to 0.185$\,GB) for
$\approx 0.1$ extra bits/weight over per-channel. Latency is flat across precisions (fake-quant);
the $\sim$15\,tok/s eager rate (clocks unpinned) is a conservative planner-only figure and excludes
perception. Realizing the memory \emph{and} a speedup on-device requires a packed-INT4 kernel
(e.g.\ Marlin/AWQ-style) and a TensorRT path---left as deployment engineering.

\section{Edge feasibility}
\begin{table*}[t]\centering
\caption{Deployment feasibility. Orin AGX row is now \emph{measured} (planner); perception
throughput remains reasoned from device class (estimates, marked $^\dagger$).}
\begin{tabular}{lccc}
\toprule
Target & memory & planner & perception (real-time) \\
\midrule
A100 (reference)      & yes & trivial & yes \\
Jetson AGX Orin 64\,GB & yes (meas.) & $2.88$\,GB, $14.9$ tok/s (meas.) & no FP32; needs TRT-INT8 port$^\dagger$ \\
Jetson Orin NX 16\,GB & workable & trivial & hard$^\dagger$ \\
Jetson Orin Nano 8\,GB & tight & ok & infeasible$^\dagger$ \\
Raspberry Pi 5 (CPU) & planner fits & sec/plan via INT4 llama.cpp$^\dagger$ & infeasible (no CUDA, custom ops)$^\dagger$ \\
\bottomrule
\end{tabular}
\end{table*}
On Orin AGX the model \emph{fits} comfortably; the obstacle is perception \emph{throughput}.
UniAD's six-camera ResNet-101 + deformable-attention BEVFormer runs FP32 with custom CUDA kernels;
reaching the $\sim$10\,Hz a vehicle needs requires exporting perception to TensorRT-INT8 (a real
engineering effort: ARM64/JetPack rebuilds of mmcv/mmdet3d ops, deformable-attention plugins). The
planner---already $<1$\,GB and INT4-ready in group-wise form---is not the constraint. On a Raspberry
Pi the gap is categorical: no CUDA and no CPU path for the deformable-attention ops, so perception
would need to be \emph{replaced}, not ported.

\section{Implication}
Quantizing the VLA language head buys memory on a component that was never the edge bottleneck, and
leaves end-to-end latency essentially unchanged because FP32 perception dominates wall-clock. The
contribution of group-wise INT4 is robustness (coherence under domain shift), not latency. The
edge-deployment research frontier for driving VLAs is \emph{perception} compression/architecture,
with the planner as a solved sub-problem.

\section{Limitations}
Footprints are measured on an A100 inference node (relative component split transfers; absolute Orin
numbers will differ). Throughput rows are device-class estimates ($^\dagger$), not on-Orin
measurements---an on-device benchmark is future work. INT4 weight footprints assume packed storage;
our research pipeline uses fake-quant (dequantized weights), so it does not yet realize the packed
memory saving.

\section{Conclusion}
For edge driving-VLAs: the planner is free (group-wise INT4, $\approx 0.26$\,GB, robust), and
perception is the wall (FP32, custom-op, latency-dominant). Optimize perception.

\begin{thebibliography}{9}
\bibitem{uniad} Hu et al. \emph{Planning-oriented Autonomous Driving (UniAD)}. CVPR 2023.
\bibitem{trtllm} NVIDIA. \emph{TensorRT-LLM / Orin INT4 deployment toolchain}. (Replace with verified citation.)
\bibitem{opendrivevla} OpenDriveVLA. 2025. (Replace with verified citation.)
\end{thebibliography}
\end{document}
